% Worked Examples: Concept 2.6.3 - Stability Margins
\documentclass[11pt,a4paper]{article}
\usepackage{worked-examples-template}

\title{\textbf{Worked Examples}\\[0.5em]
\large Concept 2.6.3: Stability Margins\\[0.3em]
\normalsize \coursesubtitle{} -- \coursetitle}
\author{Assoc.\ Prof.\ William Robertson \and Dr Sean McGowan}
\date{\today}

\begin{document}

\maketitle
\tableofcontents
\newpage

\section{Concept 2.6.3: Stability Margins}

\subsection{Example 1: Calculating Gain and Phase Margins}

\paragraph{Problem:} For the open-loop transfer function:
\begin{align}
L(s) = \frac{10}{s(s+1)(s+5)}
\end{align}

Calculate: (a) the gain crossover frequency $\omega_{gc}$ and phase margin, and (b) the phase crossover frequency $\omega_{pc}$ and gain margin.

\paragraph{Solution:}

\textbf{Part (a): Gain Crossover Frequency and Phase Margin}

The gain crossover frequency $\omega_{gc}$ is where $|L(j\omega_{gc})| = 1$ (or 0 dB).

\begin{align}
|L(j\omega)| = \frac{10}{|\omega|\sqrt{1+\omega^2}\sqrt{25+\omega^2}}
\end{align}

Setting $|L(j\omega_{gc})| = 1$:
\begin{align}
\frac{10}{\omega_{gc}\sqrt{1+\omega_{gc}^2}\sqrt{25+\omega_{gc}^2}} = 1
\end{align}

This is a transcendental equation typically solved numerically (e.g., using bisection or Newton's method). We can find an approximate solution by trial:

Starting with an initial guess based on the DC gain: at low frequency $\omega \ll 1$, $|L(j\omega)| \approx \frac{10}{\omega \cdot 1 \cdot 5} = \frac{2}{\omega}$, suggesting $\omega_{gc} \approx 2$ rad/s as a rough starting point.

Refining by iteration, try $\omega_{gc} \approx 1.4$ rad/s:
\begin{align}
|L(j1.4)| = \frac{10}{1.4\sqrt{1+1.96}\sqrt{25+1.96}} = \frac{10}{1.4 \times 1.72 \times 5.19} \approx 0.8
\end{align}

Try $\omega_{gc} \approx 1.2$ rad/s:
\begin{align}
|L(j1.2)| = \frac{10}{1.2\sqrt{1+1.44}\sqrt{25+1.44}} = \frac{10}{1.2 \times 1.56 \times 5.14} \approx 1.04 \approx 1
\end{align}

So $\omega_{gc} \approx 1.2$ rad/s.

The phase margin is:
\begin{align}
PM = 180° + \angle L(j\omega_{gc})
\end{align}

where:
\begin{align}
\angle L(j\omega) = -90° - \arctan(\omega) - \arctan(\omega/5)
\end{align}

At $\omega = 1.2$:
\begin{align}
\angle L(j1.2) &= -90° - \arctan(1.2) - \arctan(0.24) \\
&= -90° - 50.2° - 13.5° \\
&= -153.7°
\end{align}

Therefore:
\begin{align}
PM = 180° - 153.7° = 26.3°
\end{align}

\textbf{Part (b): Phase Crossover Frequency and Gain Margin}

The phase crossover frequency $\omega_{pc}$ is where $\angle L(j\omega_{pc}) = -180°$.

\begin{align}
-90° - \arctan(\omega_{pc}) - \arctan(\omega_{pc}/5) = -180°
\end{align}

Simplifying:
\begin{align}
\arctan(\omega_{pc}) + \arctan(\omega_{pc}/5) = 90°
\end{align}

Using the tangent addition formula: $\tan(A + B) = \frac{\tan A + \tan B}{1 - \tan A \tan B}$

Let $A = \arctan(\omega_{pc})$ and $B = \arctan(\omega_{pc}/5)$:
\begin{align}
\tan(90°) = \infty = \frac{\omega_{pc} + \omega_{pc}/5}{1 - \omega_{pc} \cdot \omega_{pc}/5}
\end{align}

This requires the denominator to be zero:
\begin{align}
1 - \frac{\omega_{pc}^2}{5} = 0 \quad \Rightarrow \quad \omega_{pc} = \sqrt{5} \approx 2.24 \text{ rad/s}
\end{align}

The gain margin is:
\begin{align}
GM = \frac{1}{|L(j\omega_{pc})|} \text{ (linear)} \quad \text{or} \quad GM_{dB} = -20\log_{10}|L(j\omega_{pc})|
\end{align}

At $\omega_{pc} = \sqrt{5}$:
\begin{align}
|L(j\sqrt{5})| = \frac{10}{\sqrt{5}\sqrt{1+5}\sqrt{25+5}} = \frac{10}{\sqrt{5}\sqrt{6}\sqrt{30}} = \frac{10}{\sqrt{900}} = \frac{10}{30} = \frac{1}{3}
\end{align}

Therefore:
\begin{align}
GM = \frac{1}{1/3} = 3 \text{ (linear)} = 20\log_{10}(3) \approx 9.5 \text{ dB}
\end{align}

\textbf{Results:}
\begin{itemize}
\item Gain crossover frequency: $\omega_{gc} \approx 1.2$ rad/s
\item Phase margin: $PM \approx 26°$ (marginally stable)
\item Phase crossover frequency: $\omega_{pc} = \sqrt{5} \approx 2.24$ rad/s
\item Gain margin: $GM = 3$ (or 9.5 dB)
\end{itemize}
\end{document}
